\section{Introduction}\label{intro:section}

Symplectic ball packing asks how much of a volume-preserving embedding
problem is constrained by symplectic rigidity. For \(R>0\), write
\[
 B^{2n}(R)=
 \left\{z\in\C^n:\pi\sum_{j=1}^n|z_j|^2\le R\right\},
 \qquad
 \omega_0=\sum_{j=1}^n\dd x_j\wedge\dd y_j.
\]
Thus \(R\) is the capacity, the Euclidean radius is \(\sqrt{R/\pi}\),
and \(\vol B^{2n}(R)=R^n/n!\). We ask when finitely many such balls
can be embedded symplectically into the interior of a target ball.
An embedding of a closed ball is understood to be a symplectic embedding
defined on an open neighborhood of it. For a finite disjoint union, the
compact images of the closed balls must be pairwise disjoint.

\begin{theorem}\label{intro:main}
Let \(n\ge3\) and \(k\ge1\) be integers. For positive real numbers
\(R,R_1,\ldots,R_k\), there exists a symplectic embedding
\[
 \bigsqcup_{i=1}^k B^{2n}(R_i)
 \hookrightarrow \Int B^{2n}(R)
\]
if and only if
\begin{equation}\label{intro:inequalities}
 \sum_{i=1}^k R_i^n<R^n,
 \qquad
 R_i+R_j<R\quad (1\le i<j\le k).
\end{equation}
\end{theorem}

The first inequality is the volume obstruction. The second is Gromov's
two-ball obstruction~\cite{Gromov1985}; it expresses a restriction
invisible to volume alone. Theorem~\ref{intro:main} proves that these
obstructions are complete in every real dimension at least six, resolving
Siegel and Yao's Conjecture~A positively~\cite{SiegelYao2025}. For one
ball the pairwise condition is vacuous. The strict inequalities reflect
the closed-source, open-target convention.

\subsection{Packing rigidity and stability}

Gromov's pseudoholomorphic-curve method gives the two-ball obstruction
in every dimension, and further packing obstructions in dimension
four~\cite[Section~0.3.B]{Gromov1985}. McDuff and Polterovich developed
the connection with algebraic geometry~\cite{McDuffPolterovich}.
Explicit constructions form a complementary strand: Traynor developed
symplectic packing constructions~\cite{Traynor}, and Schlenk developed
planar constructions and higher-dimensional folding
methods~\cite{SchlenkHand,SchlenkBook}. In real dimension four,
additional obstructions remain essential, so the two inequalities in
Theorem~\ref{intro:main} do not give a criterion there.

Biran's work~\cite{Biran1997,Biran1999} established packing stability:
on a closed symplectic four-manifold with rational symplectic class,
every sufficiently large number of equal balls can fill an arbitrarily
large fraction of the volume. Buse and Hind extended equal-ball
stability to higher-dimensional balls and projective spaces and then to all closed
rational symplectic manifolds~\cite{BuseHind2011,BuseHind2013}.
For unequal balls, Buse, Hind, and Opshtein proved strong packing
stability on every closed symplectic four-manifold, with all capacities
required to be sufficiently small~\cite{BuseHindOpshtein2016}.
These results identify regimes in which volume suffices; the question
here concerns every finite collection of arbitrary capacities in a ball.

Siegel and Yao formulate this higher-dimensional question and also study
stabilized packing problems~\cite{SiegelYao2025}. Their theorems on
products of four-dimensional balls with a fixed closed symplectic surface
retain four-dimensional packing obstructions. Theorem~\ref{intro:main}
concerns balls themselves, with arbitrary unequal capacities and
arbitrary finite numbers of components.

\subsection{From a mean area inequality to ball embeddings}

Normalize the target capacity to one. The common construction uses a
surface of positive genus as a base and a toric domain in \(\C^m\),
where \(m=n-1\), as a fiber. The fiber is specified by a downward
polytope \(\Delta\subset\R_{\ge0}^m\) in the moment coordinates
\(p_j=\pi|z_j|^2\): lowering any coordinate keeps a point in
\(\Delta\). Write \(H(p)>0\) for the limiting total area of the
base at fixed moment \(p\), obtained by exhausting a punctured
surface by compact bordered pieces. A ball of auxiliary capacity
\(r_i\), chosen strictly larger than the corresponding requested
capacity, requires the area profile
\[
\left(r_i-\sum_{j=1}^m p_j\right)_+,
 \qquad x_+=\max(x,0).
\]
The models contain each auxiliary simplex \(\sum_jp_j\le r_i\)
strictly away from their noncoordinate boundary faces.
The volume inequality becomes positivity of the integral of the
unused-area profile
\[
 P(p)=H(p)-\sum_i\left(r_i-\sum_jp_j\right)_+
\]
over \(\Delta\). This integral condition need not make \(P\) positive
at each moment, which is the obstacle to simply stacking the balls.

Section~\ref{cmp:section} overcomes that obstacle by Hamiltonian
comparison. Under concavity of \(P\) and positivity outside a radial
core of \(\Delta\), Lemma~\ref{cmp:comparison} constructs a smooth
toric function \(K\) and a Hamiltonian map \(q\) such that
\[
 (K-P)\circ q<K.
\]
Coordinatewise increasing rearrangements of convex functions preserve
their integrals and lead to a contraction whose fixed point gives a
uniform gap. Smooth planar disk maps realize the rearrangements with
errors smaller than that gap. Near each noncoordinate boundary face,
the generating Hamiltonians commute with its defining moment function;
this is the boundary control needed in the next step.

Section~\ref{surf:section} converts the comparison into embeddings.
A handle provides two transverse annuli and a closed one-form with a
nonzero period. These allow the Hamiltonian change of fiber coordinates
to be inserted into an exact symplectic deformation. The uniform gap
leaves disjoint layers whose area profiles contain the requested balls.
The relative Moser argument in Lemma~\ref{surf:packing} returns the
packing to the original model, preserving its boundary conditions.
The construction uses the limiting concave function \(P\); the
integrated areas of finite bordered models need only converge uniformly.

The connection between packing and Hamiltonian-group structure described
in Edtmair's abstracts~\cite{EdtmairPerfectness2025,EdtmairPacking2025}
was one motivation for the commutator viewpoint in this construction.
The present one-handle comparison and exact deformation are constructed
locally, rather than obtained from those results.

\subsection{The geometric models and the three cases}

To apply this mechanism we need models carrying nearly all the available
volume over a surface with a handle. The use of a polarization to expose
standard symplectic regions has precedents in Biran's decomposition
method~\cite{Biran2001} and Opshtein's singular polarizations and
ellipsoid packings~\cite{Opshtein2013}. Witt Nystr\"om showed that modified
deformations to the normal cone can place arbitrarily nearly all the
volume of a compact K\"ahler manifold in a normal-bundle
component~\cite{WittNystrom2024}. Here explicit toric models over
positive-genus curves supply the area profiles needed by the surface
packing lemma. Section~\ref{geo:devices} develops the relative transfer
statements that take compact packings from these models back to the
target. Its deformation and reduction tools are Moser's
method~\cite{Moser1965}, symplectic cuts~\cite{Lerman1995}, and their
K\"ahler interpretation~\cite{BurnsGuilleminLerman2002}; the relative and
equivariant details are proved where needed.

The pairwise inequalities leave three cases, summarized in
Figure~\ref{fig:proof-flow}. If every requested capacity is below
\(1/2\), Section~\ref{app:applications} degenerates \(\PP^n\)
along a complete intersection of quadrics. Its curve has positive genus
for every \(n\ge3\), and its normal model approaches the full volume
of the target. Proposition~\ref{app:small} applies the surface packing
lemma to this model.

If a requested capacity is at least \(1/2\), it is the unique such
capacity. Reserve a slightly larger central ball and pack the remaining
balls into its exterior shell. For \(n\ge4\), a hyperplane degeneration
followed by a curve degeneration in its base gives a positive-genus model
of that shell. This is Proposition~\ref{app:large}, also in
Section~\ref{app:applications}.

In complex dimension three, the same procedure would use a plane conic,
which has genus zero. Section~\ref{six:section} instead starts from a
quadric through the blowup point and uses a further degeneration along
a plane cubic with two marked points. If the reserved capacity is
\(a>1/2\), set \(s=1-a\) and \(b=(2-a)/3\). The resulting
integrated base area is
\[
 H_0(p)=9(b-p_1-p_2)-2(s-p_1-p_2)_+.
\]
Each marked point contributes one of the two subtracted logarithmic
masses. This concave density gives the required shell volume, while
exact preservation of the first circle moment \(p_1\) under transport
permits a lower
K\"ahler cut. The cubic family also supplies the invariant K\"ahler
form needed for that cut. Proposition~\ref{six:large} completes this
last case, and Section~\ref{app:completion} assembles the theorem.

All requested capacities remain arbitrary positive real numbers.
Rational auxiliary polarizations and compact truncations are chosen
with strict slack. Every transfer is defined on a neighborhood of the
relevant compact set, so the construction retains embeddings of the
whole closed source balls under finite composition.

\begin{figure}[tbp]
\centering
\begin{tikzpicture}[
  font=\small,
  box/.style={draw,align=center,inner sep=5pt,text width=39mm,minimum height=22mm},
  wide/.style={draw,align=center,inner sep=5pt,text width=118mm},
  arr/.style={-{Stealth[length=1.6mm]},semithick}]
 \node[box] (small) {All capacities \(<1/2\)\\
   Complete intersection\\of quadrics\\Proposition~\ref{app:small}};
 \node[box,right=3mm of small] (large) {One capacity \(\ge1/2\), \(n\ge4\)\\
   Two normal models\\Proposition~\ref{app:large}};
 \node[box,right=3mm of large] (six) {One capacity \(\ge1/2\), \(n=3\)\\
   Cubic with two marks\\Proposition~\ref{six:large}};
 \node[wide,below=9mm of large] (model) {Toric models over a surface with a handle:\\
   limiting unused-area profile \(P\): concave, positive mean,\\
   positive outside a radial core};
 \node[wide,below=7mm of model] (packing) {Hamiltonian comparison
   \(\longrightarrow\) disjoint ball layers\\
   Lemmas~\ref{cmp:comparison} and~\ref{surf:packing}};
 \node[wide,below=7mm of packing] (target) {Symplectic transfer to the target ball or shell;\\
   add the central ball in the two large-ball cases};
 \draw[arr] (small.south) -- ++(0,-4mm) -- ([xshift=-43mm]model.north);
 \draw[arr] (large) -- (model);
 \draw[arr] (six.south) -- ++(0,-4mm) -- ([xshift=43mm]model.north);
 \draw[arr] (model) -- (packing);
 \draw[arr] (packing) -- (target);
\end{tikzpicture}
\caption{The three geometric applications feed the same comparison and
surface-packing mechanism. Capacities are normalized by the target
capacity. The six-dimensional transfer also uses the lower K\"ahler cut.
The one-ball case is immediate.}
\label{fig:proof-flow}
\end{figure}


\subsection{Conventions and the easy direction}

We use the Hamiltonian convention
\(\iota_{X_H}\omega=-\dd H\). On \(\C^m\), moment and angle coordinates
are
\[
 p_j=\pi|z_j|^2,\qquad
 \theta_j=\frac{\arg z_j}{2\pi},\qquad
 \omega_0=\sum_j\dd p_j\wedge\dd\theta_j.
\]
Moment formulas are always interpreted in the original smooth complex
coordinates at a coordinate axis. A function is \emph{toric} if it
depends only on the moments. Projective spaces and line bundles use
Chern class units: a projective line in the unit class has area one.
Projectivizations parametrize lines, and \(U=c_1(\cO(1))\) denotes the
dual tautological class. We normalize \(\dd^c\) by
\(\dd^c\log|z|^2=\dd\theta\) away from zero; consequently
\(\dd\dd^c\log|z|^2\) has a unit atom at zero.

\begin{proof}[Necessity in Theorem~\ref{intro:main}]
The union of the finitely many compact images is contained in
\(\Int B^{2n}(\sigma)\) for some \(\sigma<R\).
Volume preservation gives
\(\sum_iR_i^n\le\sigma^n<R^n\). Gromov's two-ball obstruction
\cite[Section~0.3.B]{Gromov1985}, in the capacity convention
of~\cite[Theorem~1.1]{SiegelYao2025}, gives
\(R_i+R_j\le\sigma<R\) for every pair.
\end{proof}

Conjugating the source and target by the same dilation reduces
sufficiency to target capacity one; the conformal factors in the two
pullbacks cancel. We henceforth use this normalization. A single ball
of capacity less than one embeds by inclusion. The proof for multiple
balls is completed after the three geometric applications, in
Section~\ref{app:completion}.
